\let\negmedspace\undefined
\let\negthickspace\undefined
\documentclass[journal]{IEEEtran}
\usepackage[a5paper, margin=10mm, onecolumn]{geometry}
%\usepackage{lmodern} % Ensure lmodern is loaded for pdflatex
\usepackage{tfrupee} % Include tfrupee package

\setlength{\headheight}{1cm} % Set the height of the header box
\setlength{\headsep}{0mm}     % Set the distance between the header box and the top of the text

\usepackage{gvv-book}
\usepackage{gvv}
\usepackage{cite}
\usepackage{amsmath,amssymb,amsfonts,amsthm}
\usepackage{algorithmic}
\usepackage{graphicx}
\usepackage{textcomp}
\usepackage{xcolor}
\usepackage{txfonts}
\usepackage{listings}
\usepackage{enumitem}
\usepackage{mathtools}
\usepackage{gensymb}
\usepackage{comment}
\usepackage[breaklinks=true]{hyperref}
\usepackage{tkz-euclide} 
\usepackage{listings}
% \usepackage{gvv}                                        
\def\inputGnumericTable{}                                 
\usepackage[latin1]{inputenc}                                
\usepackage{color}                                            
\usepackage{array}                                            
\usepackage{longtable}                                       
\usepackage{calc}                                             
\usepackage{multirow}                                         
\usepackage{hhline}                                           
\usepackage{ifthen}                                           
\usepackage{lscape}
\begin{document}

\bibliographystyle{IEEEtran}
\vspace{3cm}

\title{9.4.46}
\author{EE25BTECH11033 - Kavin}
% \maketitle
% \newpage
% \bigskip
{\let\newpage\relax\maketitle}

\renewcommand{\thefigure}{\theenumi}
\renewcommand{\thetable}{\theenumi}
\setlength{\intextsep}{10pt} % Space between text and floats
\textbf{Question}:\\
A pole has to be erected at a point on the boundary of a circular park of diameter $1.3$ metres in such a way that the difference of its distances from two diametrically opposite fixed gates A and B on the boundary is $7$ metres. Is it possible to do so? If yes, at what distances from the two gates should the pole be erected?\\    
\bigskip

\textbf{Solution}:\\

Let the circular park have gates $\vec{A}$ and $\vec{B}$ at the ends of a diameter. The distance between them, $AB$, is the diameter, which is $13$ m. Let the pole $\vec{P}$ be a point on the boundary of the circle.

\begin{itemize}
    \item The angle subtended by the diameter at any point on the circumference is a right angle ($90\degree$). Therefore, $\triangle APB$ is a right-angled triangle with the hypotenuse being the diameter AB.
\end{itemize}


Using Baudhayana's theorem for $\triangle APB$
\begin{align*}
 PA^2 + PB^2 = AB^2     
\end{align*}


\begin{align}
PA^2 + PB^2 = 169
\end{align}

We are also given that the difference of the distances from the pole to the gates is 7 meters:
\begin{align}
\abs{PA - PB} = 7    
\end{align} 
Assuming $PA > PB$, we have:
\begin{align}
PA - PB = 7
\end{align}
\begin{align}
\implies PB = PA - 7
\end{align}
\begin{align}
\implies PA^2 + \brak{PA - 7}^2 = 169
\end{align}
Let $PA$ be $x$,
\begin{align}
\implies x^2 + \brak{x - 7}^2 = 169
\end{align}
\begin{align}
\implies    y = x^2 - 7x - 60 = 0
\end{align}
which 
 can be expressed as the conic
\begin{align}
    \vec{x}^\top\vec{V}\vec{x} + 2\vec{u}^\top\vec{x} + f = 0 \\
    \vec{V} = \myvec{1 & 0 \\ 0 & 0}, \vec{u} = \myvec{-\frac{7}{2} \\ -\frac{1}{2}}, f = -60.
\end{align}
To find the roots of the expression, we find the points of intersection of the conic with the $x$-axis 
\begin{align}
    \vec{x} = \vec{h} + k\vec{m} \\
    \vec{h} = \myvec{0 \\ 0}, \vec{m} = \myvec{1 \\ 0}
\end{align}\\
\bigskip 

The points of intersection of the line 
\begin{align}
L: \quad \vec{x} = \vec{h} + \kappa \vec{m} \quad \kappa \in \mathbb{R}
\label{eq:conic_tangent}
\end{align}
with the conic section is given by
\begin{align}
\vec{x}_i = \vec{h} + \kappa_i \vec{m}
	\label{eq:chord-pts}
\end{align}
%
where
\begin{multline}
\kappa_i = \frac{1}
{
\vec{m}^{\top}\vec{V}\vec{m}
}
\lbrak{-\vec{m}^{\top}\brak{\vec{V}\vec{h}+\vec{u}}}
%\\
\pm
%{\small
\rbrak{\sqrt{
\sbrak{
\vec{m}^{\top}\brak{\vec{V}\vec{h}+\vec{u}}
}^2
	-\text{g}
\brak
{\vec{h}
%\vec{h}^{\top}\vec{V}\vec{h} + 2\vec{u}^{\top}\vec{h} +f
}
\brak{\vec{m}^{\top}\vec{V}\vec{m}}
}
}
%}
\label{eq:tangent_roots}
\end{multline}
The values of $k$ are given by
\begin{align}
k_i    &= \frac{1}{1} \brak{\frac{7}{2} \pm \sqrt{\brak{\frac{7}{2}}^2 + 60}} \\
  \implies  k_1 &= -5, k_1 = 12
\end{align}
Hence the points of intersection are
\begin{align}
    \vec{h} + k\vec{m} = \myvec{-5 \\ 0}, \myvec{12 \\ 0}
\end{align}

Hence the solutions 
are $x = -5$ and $x = 12$. We reject $x = -5$ as the length of the side cannot be negative.\\
This gives the value $PA = 12$ and $PB = 5$.\\
\bigskip

Therefore, it is possible to erect the pole as described.
\newpage
\begin{figure}[H]
    \centering
    \includegraphics[width=0.6\columnwidth]{figs/fig1.png}
    \caption{Intersection of $y = x^2 - 7x - 60$ with the $x$-axis}
	\label{fig:quad1}
\end{figure}
\begin{figure}[H]
    \centering
    \includegraphics[width=0.6\columnwidth]{figs/fig2.png}
	\label{fig:quad2}
\end{figure}
\end{document}


